\chapter{API Documentation}
\label{ch:api-documentation}

The system exposes three classes of interface:

\begin{enumerate}
  \item the \textbf{registry HTTP interface} (read-only static JSON);
  \item the \textbf{programmatic interfaces} of the published packages;
  \item the \textbf{CLI interface}, documented in Chapter~\ref{ch:cli-reference}.
\end{enumerate}

There is no authenticated API, no write API and no server-side application.

\section{Registry HTTP interface}

\subsection*{Base URL resolution}

\begin{longtable}{@{}L{4.2cm}L{10cm}@{}}
\caption{Registry base URL resolution}\\
\toprule
\textbf{Consumer} & \textbf{Resolution chain} \\
\midrule
\endfirsthead
\toprule
\textbf{Consumer} & \textbf{Resolution chain} \\
\midrule
\endhead
CLI & \texttt{--registry <url>} flag $\rightarrow$ \texttt{EXPO\_SNIPS\_REGISTRY}
environment variable $\rightarrow$ empty (embedded catalog). \\
Extension & \texttt{expo-snips.registryUrl} setting $\rightarrow$ empty (embedded
catalog). \\
Declared default constant & \texttt{https://registry.expo-snips.dev}
(\texttt{DEFAULT\_REGISTRY} in both clients) --- exported but not used as an
implicit default; an empty base selects the embedded catalog. \\
Dev server & \texttt{pnpm serve} $\rightarrow$ \texttt{http://localhost:3000}. \\
\bottomrule
\end{longtable}

The base is normalised by trimming whitespace and removing trailing slashes.

\subsection*{\texttt{GET /index.json}}

\begin{longtable}{@{}L{3.4cm}L{10.8cm}@{}}
\toprule
\textbf{Aspect} & \textbf{Detail} \\
\midrule
\textbf{Purpose} & Return the catalog index used for browsing and selection. \\
\textbf{Method} & GET \\
\textbf{Authentication} & None. \\
\textbf{Authorization} & None (public read). \\
\textbf{Headers} & None required. \\
\textbf{Parameters} & None. \\
\textbf{Request body} & None. \\
\textbf{Response} & JSON array of index entries. \\
\textbf{Timeout} & Client-side abort after 10\,000\,ms. \\
\textbf{Status codes} & Any non-2xx is treated as failure by the client
(\texttt{response.ok} check); the client then falls back to the embedded catalog. \\
\textbf{Validation} & The client accepts the payload only when it is an array with
at least one element. \\
\textbf{Source} & Produced by \texttt{scripts/build-registry.mjs} via
\texttt{toIndexEntry}. \\
\bottomrule
\end{longtable}

Index entry shape:

\begin{lstlisting}[language=JavaScript]
interface RecipeIndexEntry {
  id: string;          // e.g. "auth/login"
  name: string;        // e.g. "Login Screen"
  description: string;
  version: string;     // semver, e.g. "0.3.0"
  typescript: boolean;
  tags: string[];
  supported: {
    styling: string[];
    router: string[];
    ui: string[];
    icons: string[];
  };
}
\end{lstlisting}

Example response fragment (from \texttt{registry/index.json}):

\begin{lstlisting}[language=JavaScript]
[
  {
    "id": "app/home",
    "name": "Home Dashboard Screen",
    "description": "Home screen with greeting, stats cards and upcoming tasks list.",
    "version": "0.1.0",
    "typescript": true,
    "tags": ["app", "home", "dashboard"],
    "supported": {
      "styling": ["nativewind", "stylesheet"],
      "router": ["expo-router"],
      "ui": ["native"],
      "icons": ["lucide", "none"]
    }
  }
]
\end{lstlisting}

\textbf{Note.} The index omits \texttt{variants}, so clients browsing only the
index cannot tell whether a recipe offers variants. Only the full recipe document
carries them.

\subsection*{\texttt{GET /recipes/\{id\}.json}}

\begin{longtable}{@{}L{3.4cm}L{10.8cm}@{}}
\toprule
\textbf{Aspect} & \textbf{Detail} \\
\midrule
\textbf{Purpose} & Return one complete recipe with templates inlined. \\
\textbf{Method} & GET \\
\textbf{Path parameter} & Recipe id, URL-encoded by the client
(\texttt{encodeURIComponent}); note that the on-disk path uses a slash
(\texttt{recipes/auth/login.json}) while the client encodes the slash as
\texttt{\%2F}. \\
\textbf{Authentication} & None. \\
\textbf{Response} & A \texttt{ScreenRecipe} object with \texttt{files[].template}
strings. \\
\textbf{Client acceptance} & Accepted only when the payload is an object with a
truthy \texttt{id}. \\
\textbf{Fallback} & Embedded catalog lookup by id. \\
\textbf{Source} & Written by \texttt{scripts/build-registry.mjs} to
\texttt{registry/recipes/<id>.json}. \\
\bottomrule
\end{longtable}

Example (abridged) response:

\begin{lstlisting}[language=JavaScript]
{
  "id": "auth/login",
  "name": "Login Screen",
  "version": "0.3.0",
  "framework": "react-native",
  "runtime": "expo",
  "typescript": true,
  "supported": {
    "styling": ["nativewind", "stylesheet"],
    "router": ["expo-router", "react-navigation"],
    "ui": ["native", "shadcn"],
    "icons": ["lucide", "none"]
  },
  "dependencies": ["react-hook-form", "zod", "@hookform/resolvers", "lucide-react-native"],
  "files": [
    {
      "path": "app/login/index.tsx",
      "when": { "router": ["expo-router"], "ui": ["native"] },
      "variants": { "glass": "import { ... } from 'react-native'; ..." },
      "template": "import { View, Text, ... } from 'react-native'; ..."
    }
  ],
  "variants": [{ "id": "glass", "label": "Glassmorphism", "description": "..." }],
  "tags": ["auth", "login", "form"],
  "primitives": ["ui/utils", "ui/button", "ui/input"]
}
\end{lstlisting}

\subsection*{Local development server}

\texttt{scripts/serve.mjs} starts a Node \texttt{http} server that serves the built
\texttt{registry/} folder as static files (default port 3000, overridable with the
\texttt{PORT} environment variable). It:

\begin{itemize}
  \item answers \texttt{OPTIONS} with \texttt{204} and permissive CORS headers
        (\texttt{Access-Control-Allow-Origin: *}, methods \texttt{GET, OPTIONS});
  \item resolves the URL path against the registry root and rejects any path that
        escapes it, preventing traversal;
  \item maps directories and trailing slashes to \texttt{index.json};
  \item returns \texttt{404} with \texttt{\{"error":"Not found"\}} for unknown paths;
  \item sets a content type from a small extension map (JSON, HBS, PNG, JPG, GIF,
        SVG) and otherwise \texttt{application/octet-stream}.
\end{itemize}

It refuses to start when \texttt{registry/index.json} is missing and instructs the
operator to run \texttt{pnpm build:registry} first.

\section{Programmatic interfaces}

\subsection*{\texttt{@expo-snips/core} (MIT, published)}

\begin{longtable}{@{}L{6.4cm}L{7.8cm}@{}}
\caption{Public exports of the core package}\\
\toprule
\textbf{Export} & \textbf{Purpose} \\
\midrule
\endfirsthead
\toprule
\textbf{Export} & \textbf{Purpose} \\
\midrule
\endhead
\texttt{renderTemplate(templateStr, context)} & Render one template to a string. \\
\texttt{getFilesForRecipe(recipe, context, resolveFile?)} & Render a recipe into
\texttt{\{path, content\}[]}. \\
\texttt{GenerationContext}, \texttt{GenerationOptions}, \texttt{RenderContext} & Input
types for rendering. \\
\texttt{ScreenRecipe}, \texttt{FileTemplate}, \texttt{ScreenVariant},
\texttt{Conditions} & Data model types. \\
\texttt{StylingOption}, \texttt{RouterOption}, \texttt{UIOption}, \texttt{IconOption} &
Option unions. \\
\bottomrule
\end{longtable}

Usage example:

\begin{lstlisting}[language=JavaScript]
import { getFilesForRecipe, GenerationContext, ScreenRecipe } from "@expo-snips/core";

const context: GenerationContext = {
  projectRoot: process.cwd(),
  options: { styling: "nativewind", router: "expo-router", ui: "native", icons: "lucide", typescript: true },
};

const files = getFilesForRecipe(recipe as ScreenRecipe, context);
// => [{ path: "app/login/index.tsx", content: "..." }, ...]
\end{lstlisting}

\subsection*{\texttt{@expo-snips/cli} (MIT, published)}

The CLI package exposes no documented programmatic API. Its
\texttt{main} points to \texttt{dist/index.js}, which executes the commander
program as a side effect when imported, so it is intended for use as a binary
rather than as a library.

\subsection*{\texttt{@expo-snips/registry} (private)}

\begin{longtable}{@{}L{6.4cm}L{7.8cm}@{}}
\caption{Exports of the registry package}\\
\toprule
\textbf{Export} & \textbf{Purpose} \\
\midrule
\endfirsthead
\toprule
\textbf{Export} & \textbf{Purpose} \\
\midrule
\endhead
\texttt{loadRecipeFromFile(filePath)} & Parse a recipe and inline \texttt{templateFile},
variant and per-styling references. \\
\texttt{getAllRecipes()} & Recursively load and inline every recipe under
\texttt{data/}. \\
\texttt{getRecipeById(id)} & Load and match a single recipe. \\
\texttt{searchRecipes(query)} & Case-insensitive match on id, name or tags. \\
\texttt{fetchRecipes(url)} & Fetch a remote registry document; accepts a bare array,
a \texttt{\{recipes\}} bundle, or a single recipe object; throws on a non-OK
response or an unrecognised body. \\
\bottomrule
\end{longtable}

Its filesystem functions resolve paths relative to the package directory using
\texttt{\_\_dirname}, so they are only meaningful inside the repository or a
distribution that includes \texttt{data/} and \texttt{templates/}.

\subsection*{\texttt{expo-snippets} (VS Code extension, proprietary)}

Contributes one command and registers a second at runtime.

\begin{longtable}{@{}L{5.6cm}L{8.6cm}@{}}
\caption{Extension commands and settings contributions}\\
\toprule
\textbf{Interface} & \textbf{Definition} \\
\midrule
\endfirsthead
\toprule
\textbf{Interface} & \textbf{Definition} \\
\midrule
\endhead
Command & \texttt{expo-snips.addScreen} --- title ``Expo Snips: Add Screen...''. \\
Command (runtime) & \texttt{expo-snips.generateFromCompletion(recipeId?)}. \\
Setting & \texttt{expo-snips.registryUrl} (string, default empty). \\
Setting & \texttt{expo-snips.defaultStyling} (enum, default \texttt{nativewind}). \\
Setting & \texttt{expo-snips.defaultRouter} (enum, default \texttt{expo-router}). \\
Setting & \texttt{expo-snips.defaultUi} (enum, default \texttt{native}). \\
Setting & \texttt{expo-snips.defaultIcons} (enum, default \texttt{lucide}). \\
Setting & \texttt{expo-snips.installDependencies} (boolean, default \texttt{true}). \\
Activation & \texttt{onCommand:expo-snips.addScreen} and
\texttt{onLanguage:\{typescript,\allowbreak typescriptreact,\allowbreak javascript,\allowbreak javascriptreact\}}. \\
\bottomrule
\end{longtable}

\section{Failure and error semantics summary}

\begin{longtable}{@{}L{4.4cm}L{3.4cm}L{6.4cm}@{}}
\caption{Error semantics across interfaces}\\
\toprule
\textbf{Failure} & \textbf{Signal} & \textbf{Client behaviour} \\
\midrule
\endfirsthead
\toprule
\textbf{Failure} & \textbf{Signal} & \textbf{Client behaviour} \\
\midrule
\endhead
Registry unreachable & Fetch rejects or aborts & Fall back to the embedded catalog;
CLI prints a warning when a base was configured. \\
Registry returns non-OK & \texttt{response.ok} false & Same fallback. \\
Registry returns empty array & Array length 0 & Same fallback. \\
Recipe missing & \texttt{getRecipeById} returns null & CLI: error and exit 1.
Extension: error message. \\
Manifest unreadable & JSON parse throws & Treated as empty manifest. \\
Unknown primitive id & Recipe lookup returns null & Warning; primitive skipped. \\
Install failure & Child process error & CLI aborts with the process error; files
are already written. \\
\bottomrule
\end{longtable}
